\section{流动与修复：人怎样把帝国的断处带到别处}
\label{sec:synthesis-mobility-repair}

迁徙不是两汉编年中的余波。它既可能起于战乱、饥荒、徭役和边地征发，也可能经由婚姻、市场、求学、仕进和服役缓慢发生。人离开原籍时带走的并不只是劳力：还有亲属关系、债务、技艺、地方记忆和对官府的经验。到达新地以后，这些东西会改变谁能耕田、谁可被编户、谁有资格请求保护。西汉的实边、新末的流民、东汉度田后的安置、汉末的屯田与区域军政，因而应当放在一条迁徙与修复的线索中理解。

\subsection{离开原籍：国家的分类与家庭的时机}

朝廷常以徙民、实边、流民、降人等名词描述移动人口。这些名称使行政能够处理人数、地点与供给，却不能说明每个人何时决定离开、与谁同行、能带走什么。一个家庭可能因河渠失修、兵役临近或欠债而先派出一名成员探路；也可能在一夜之间因武装接近而逃散。传世史书通常留下大规模迁徙的原因和结果，难以保存这种决定的细部。叙事不能为沉默补造对话，却可以把迁徙放回道路、渡口、农时和粮食这些可见的条件。

西汉的边地安置尤其显示行政目标与家庭时机并不一致。朝廷希望河西、朔方等地有人耕作、守塞和供应传舍；迁来者则要判断水源、气候、邻近部众和市场是否足以维持生活。实边可以扩大可耕地和守备，也会使新旧居民围绕水、地和劳力重新安排关系。政策文本最能证明朝廷想把谁迁到哪里，较难说明当地人如何接纳或排斥，故不能把“置郡”“徙民”写成一项已经完成的社会事实。

\subsection{新朝与复汉：流民不是没有政治判断的人}

新末的灾害、兵乱和制度改易，使离开原籍成为许多人的生存选择。把流民写成被动的人群，会遮蔽他们如何寻找粮食、亲属、城邑和保护者；把他们全写成反叛者，又会把求生误作统一的政治意志。绿林、赤眉及地方集团能够吸纳人群，部分因为它们提供了粮食、武装和可辨认的名义。不同地方的人投向不同网络，不是因为他们共享一种抽象的“民心”，而是因为道路、仓储、地方首领和战争风险不同。

复汉政权要稳定，必须使这些移动的人重新拥有可预期的位置。雒阳的任命、赦令和度田并非仅为恢复税源，也为了让安置、归附与征发重新有规则。规则会带来新的压力：谁被登记、谁能保留暂住地、谁必须回到原籍，都是家庭与地方官需要处理的问题。东汉的国家能力就在这类不显眼的协商中形成，而不只在都城的礼仪中形成。

\subsection{道路上的知识与宗教}

人随道路移动，文本与仪式也随道路移动。经师、学生、官员和书吏把经典、荐举与官署语言带往不同地区；商旅、戍卒、移民和地方供养者则可能携带治疗、祈福、佛教相关语汇和各类方术实践。这样的流动不意味着所有地方接受同一种信念。楚王英案中的供养语汇能说明某些实践进入了官方可识别的范围，不能证明东汉各地已有同样的教团或寺院。

石经、太学和地方师承让知识获得跨地区的共同参照，也让资格与声望可以离开原乡。它们在东汉中期成为政治问题，并非知识忽然脱离社会，而是仕进、评价和宫廷人事已经紧密相连。黄巾时期的宗教语言又显示，当赈济、治安和地方保护失去可信度时，治疗、共同禁忌和末世解释能够成为组织人群的方式。不能从起事倒推出所有仪式团体都是叛乱预备队；应当问哪些地方的家庭和流动者已经缺少其他可用的联结。

\subsection{边费的迁徙后果}

边防并不只消耗国库，它也重排人口。远征与驻防需要戍卒、运输者、屯田者和传舍人员；战争又可能使边民、降人和附近乡里的家庭离开原地。河西道路、羌地山谷和北边市场各有不同的安置难题。中枢以互市、屯田、赏赐或迁徙处理风险时，成本会经由粮食、劳役和土地落到具体地区。

东汉边费更加容易转化为地方军政网络。州郡要筹粮、招募、安置流民，便需要与大族、商人、部曲和地方首领合作。这样的合作可能让一地度过危机，也会使掌握仓储与武装的人获得新的谈判位置。汉末区域中心的力量，部分来自它们能够把迁徙人口、可耕地和军事供给重新接合；这不是边政之外的结果，而是边政长期成本在地方社会中的回流。

\subsection{屯田：安置、供给与新的约束}

屯田常被描述为将荒地和流民同时利用的方案。它的实际意义更复杂。对于竞争者，屯田能够稳定军粮、减少掠夺并将人口纳入可管理的编制；对于参与者，它可能提供相对稳定的耕作机会，也可能意味着新的劳动义务、监督和迁徙限制。材料不足以让我们为所有地区、所有屯田者下同一种判断，因而不应把屯田写成纯粹恩惠或纯粹压迫。

曹操的屯田之所以与奉迎献帝相连，正在于供给与名义需要彼此支持。皇帝诏令能使任命和征发具有公共语言，粮仓、道路和兵员则使这种语言在北方产生实际后果。魏、蜀、吴后来各自处理人口、土地和军粮，也说明二二〇年的受禅没有结束迁徙与修复的问题。它只是使这些问题不再由同一个中心结算。

\subsection{材料怎样看见流动，又怎样看不见}

户籍和诏令会把移动人口压成可登记的类别；简牍有时记录具体人员、凭验和粮食；墓葬、碑刻和宗教文本则让少数人的家庭与供养关系露出痕迹。每一类材料都只照见流动的一面。户籍不能代表未登记者，简牍不能代表全部道路，后出的宗教传记也不能替代早期实践本身。把它们并读，不是为了拼出一份完美的人口地图，而是为了区分国家想要安置什么、某些地方实际留下了什么，以及哪些经验仍然无法被材料捕捉。

\subsection{编户的门槛：一户人家怎样重新进入国家的视野}

把两汉的户籍理解成一张静止的人口表，最容易错过迁徙的实际困难。县廷需要一个可用的名字、居处和承担赋役的单位，才知道应向谁传达徭役、由谁领受赈贷、由谁为邻里作保；家庭却未必以同样的边界生活。战后有人死去，有人从原籍逃来，有人寄居在姻亲或强宗的土地上，有人先以雇工、佃作或服役的身份落脚。对官府而言，这些人必须在“本县”“客”“逃亡”或其他可处理的类别中取得位置；对当事人而言，任何类别都可能同时带来粮食、土地、保护和征发。故而编户不是档案室中的收尾步骤，而是人能否留在一地的谈判。

早汉的材料让我们首先看见中枢怎样把人放入空间。前一九六年前后，刘敬提出迁徙关东豪族以实关中的建议；《史记·刘敬叔孙通列传》与《汉书·高帝纪》留下这一政策的轮廓，而所称人数和实际执行范围并不相同，不能把它换算成全国人口流动的比例。可以较有把握地说，长安周边的稳定需要人户、耕地与政治控制一道重组；不能据此断言每个被迁者的家产如何清点、是否全家同行，或关东地方从此不再有可动员的精英。史料的缝隙本身很重要：国家能够宣布把谁迁往何处，却难以留下每一户如何在新县取得水源、婚姻和邻里关系的连续记录。

这也说明户籍的效力从不只来自强制。一个迁入家庭若能得到耕作机会、知道粮食从何处取得、相信来年的登记不会任意翻覆，才有理由把劳力和亲属关系固定在新地。反过来，官府若要取得稳定的税役来源，也必须允许地方通过保伍、宗族、雇佣和市场吸收一部分不规则的人口。传世纪传偏爱诏令和大臣的主张，居延、肩水金关、悬泉置等简牍则在较晚的边塞行政中保存了凭验、传递和供给的零星痕迹。两类材料不能合成一张覆盖全国的户口图，却共同提醒我们：可见的人口是通过文书、道路、担保和给养逐项做出来的。

编户还会改变家庭内部的风险分配。国家以户为征发和赈贷的单位，未必等于户内的人拥有相同的土地、粮食或离开的能力。老人、妇女、儿童、寄居者和暂时外出的劳力，在文本里常被压入一个户名之下；他们在灾年如何取得口粮，在迁移时由谁照料，通常没有直接记录。承认这一限度并非放弃社会史，反而要求叙事不把“户数增加”直接改写为每家生活安定，也不把“流民”当成没有年龄、性别和亲属关系的队列。户籍为国家提供了一个结算尺度，却不能穷尽家庭维持自身的尺度。

\subsection{仓廪与去留：赈济不是户籍之外的慈惠}

灾害到来时，户籍的边界会变得格外尖锐。谷物在仓中、道路能通、县吏知道谁在本地，这几项条件缺一不可；即使诏令宣布赈贷，粮食也必须先从某处移出，再由某些人认领、运输和分配。西汉帝纪与《汉书·食货志》保存了减租、赐谷、开仓和劝课等不同语言，它们能说明朝廷把饥荒、赋役和恢复耕作视为相连的问题，却不能证明每次命令在每个县都按同一速度到达。尤其不能把史书所记的大数当作每一户实际所得：申报的数量、输送中的损耗、地方优先次序和未登记者的处境，大多不在同一条材料中。

不过，正因为赈济需要识别对象，它才会反过来塑造户籍。地方官若只向旧簿上的人发粮，刚到的流民可能被排除；若为了止乱而承认暂住者，又会改变谁被看作本地的负担与资源。一个家庭在灾年选择留下，不一定是出于对朝廷的抽象忠诚，也可能因为亲族能借粮、市场尚有交易、县仓仍可接近；选择离开，也不必然表示反叛，可能只是种子、役期和道路风险已经无法承受。把赈济放在这些选择中，才能看出“救荒”既是维持生命的措施，也是防止一地人户和耕地一同空掉的治理技术。

新朝的危机把这一关系推到更不稳定的状态。王莽改易土地、币制和官名的政策，在《汉书·王莽传》中以强烈的成败叙事保存下来；同时期灾害、战乱和地方动员却使任何统一的登记与给养都难以稳步执行。问题不只是政策是否合乎某种理想，而是县乡能否用仍被人相信的货币和文书组织借贷、征收与赈给。新币若难以在市场兑换，土地名目若与既有占有和债务冲突，家庭便会在官府的分类、地方保护与求生需要之间另找可行路径。我们无法从后出的叙事逐户复原这种路径，但也不能用“改革失败”四字抹掉它。

因此，绿林、赤眉等人群的扩张不能只被解释为中央秩序崩塌后的情绪。饥者要知道哪里有粮，携家而行者要知道渡口、城邑和熟人是否仍可到达；一支武装若能供给、分配和维持某种内部约束，便可能吸收原本没有共同政治目标的人。史书对于这些集团的兵数、行动路线和口号常有夸张或立场，足以提示而不足以精确重建各地流民规模。更稳妥的结论是：当户籍不能连接到谷物和保护时，家庭会把自己的生计接到别的网络上；新政权若不能重新建立这种连接，国号的更替便难以成为日常秩序。

\subsection{度田之后：复兴怎样处理旧簿上没有的人}

东汉复兴后的度田，通常被当作光武朝恢复财政的节点。它当然涉及土地与赋役，但若从家庭看，问题还包括谁能以何种身份留在新处。建武时期的战乱已经改变许多地方的人地关系：旧主死亡或离散，田地荒芜或被他人耕作，流入者与原住者共用道路、市场和水源。重新丈量和登记能让朝廷估计可征发的基础，也会迫使地方把暂居、依附、隐匿和争议中的关系说成一个确定的数字。《后汉书·光武帝纪》与相关列传记载度田引起的反应，能够说明政策确曾触及地方利益；它们不能证明所有抵抗出自同一阶层，更不能让我们将每一条数字当作无误的社会普查。

这场工作需要地方中介，因而不可能只是中枢单向的“看见”。里正、乡吏、豪右、宗族长者和能够书写文书的人，知道某块田由谁实际耕种、哪一家仍有可以承担劳役的成员、哪些人可被说成客居或逃亡。正是这种知识使登记成为可能，也使它能够被选择性使用。中央想减少隐漏，地方则可能在保护亲属、维持劳力与避免额外负担之间调节口径。将度田仅写作国家胜过地方，或写作地方完全挫败国家，都会忽略它是一连串局部协商；国家并非没有力量，地方也并非不受命令约束，双方却都无法在没有对方信息的情况下完成簿籍。

赦令和赈给在此后的意义，也不应被缩小为朝廷的恩典展示。它们为战后人口提供重新归属的可预期条件：有人可因赦免不再担心旧罪，有人可因给养渡过播种前的缺粮，有人则通过服役、屯垦或依附取得暂时位置。每一种路径都包含代价。免除旧负担会削弱短期收入；把人安置在边地或屯田会把风险转移到新的水土和军政纪律中；承认地方保护又会增强保护者的议价能力。东汉之所以能够延续，并非因为这些矛盾被消除，而是因为在一些时期，文书、道路、仓储和地方协作仍足以让它们不至于同时失控。

\subsection{晚汉的客籍：地方保护怎样改变国家的结算}

到东汉末年，灾异、疫病、兵乱与反复征发使“原籍”越来越不能保证安全。黄巾起事及其后的军事竞争，在《后汉书·灵帝纪》《后汉书·皇甫嵩传》等材料中主要以朝廷平乱和将领行动呈现；这些记载能确证广泛动员和军事破坏，却不能替我们计算每一地的饥饿者或逃亡户数。更贴近家庭的提问是：谁在县城失守前先走，谁留在坞壁，谁被军队征作运输者，谁以客居身份进入大姓、寺观或地方武装的保护。答案因地区而异，不能用一个“流民社会”概括。

地方保护之所以有力量，不仅因为它拥有武器。它还可能控制仓谷、熟悉小路、能向外来军队交涉，并能为到来者提供某种可预期的名分。坞壁和宗族网络未必天然善待寄居者；接纳劳力可以扩大耕作和守卫，也可能制造债务、依附和排斥。对地方首领来说，登记新来者既意味着增加可用人手，也意味着承担粮食和安全风险；对流民来说，进入这种网络可能比无处投奔更好，却不等于恢复了自由的旧生活。国家簿籍在这里并未完全消失，而是与私人的担保、军队的名册和区域政权的征发并列，彼此争夺谁有权定义一户人家的归属。

曹操在北方推行屯田的材料，尤其《三国志·魏书·武帝纪》及后出的编年叙事，常把政策置于军粮和统一竞争的框架。这个框架是必要的，却仍需看到屯田把流动家庭放进了新的登记、生产和军事关系。它可能减少无序掠夺、给部分人以可耕之地，也可能以定额、监督和迁徙限制换取稳定供给。不同地区、不同身份的参与者不可能有相同经验，材料也不足以宣称他们普遍受益或普遍受害。可以确定的是，能够把人、地、粮和军队接合的区域中心，比只保有帝国名义的中心更有能力决定流民的去留。

二二〇年的禅代没有使这些客籍、屯田户和依附关系自动归一。魏、蜀、吴各自继承县邑、文书和征发的语言，却要在不同的道路、仓储和人口分布中处理同样的问题。汉末所谓地方碎裂，因而不是地方忽然拒绝一切制度，而是户籍、赈济和保护不再由同一套中心规则同时结算。家庭在危机中作出的去留选择，既是这种碎裂的后果，也不断为新的区域秩序提供劳力、税源和兵源。
\subsection{制度得失：修复为何总有新的裂缝}

安置、度田、实边与屯田都试图把移动的人重新接入土地、赋役和保护。它们使国家在战后或边地危机中重新获得人力与粮食，也使家庭有机会避免完全失去生计。代价在于，重新接入总要决定谁被登记、谁承担守卫、谁有资格取得土地，因而会在新地方制造新的排除与依附。

两汉不断修复这些裂缝，却不能让裂缝完全消失。西汉用行政与边地经营扩大可调度的空间；新朝显示命令若脱离地方信任会怎样失速；东汉重建了簿籍、学校与道路，又在晚期看到这些接口被区域网络分别占有。迁徙史把这种连续与断裂放在家庭和道路的尺度上，使读者看到帝国并非只在宫门中改变，也在每一次离开、到达和重新登记中改变。
